% Shared by main.tex and extended/main.tex. Lines marked % CHECK-DATA depend on the data.
\section{Introduction}
\label{sec:intro}

Developers of large language models (LLMs) run safety evaluations before a release and report them in system cards~\citep{shevlane2023extreme,anthropic2026rsp,openai2025preparedness}; a virology question set~\citep{gotting2025vct} runs in minutes through a programming interface. Physical-AI safety evaluations test a robot, a vehicle or a model that controls one, from question answering about physical danger~\citep{jindal2025danger} to world-model rollouts~\citep{geminirobotics2025veo}, jailbreaks of real robots~\citep{robey2024robopair} and vehicle track protocols~\citep{iihs2024paeb}. Many need hardware, scenes and staff, so they look too expensive to be worth running.

Cost alone cannot say whether an evaluation is worth running, because an evaluation is a measurement, worth what its result does. A safety evaluation is worth the harm its result avoids by improving the release decision, so a cheap test whose result cannot move the decision is worth nothing and an expensive test that would stop a harmful release is worth much. Decision analysis quantifies this as the expected value of sample information (EVSI)~\citep{raiffa1961applied,howard1966information}, used in health economics to choose which studies to fund~\citep{claxton1999irrelevance}. AI safety evaluations deserve the same scrutiny: their budget is finite and should go to the evaluations that give the most safety per dollar.

We estimate the value of information of \voiNPhysical{} physical-AI safety evaluations and compare it with that of \voiNLLM{} LLM safety evaluations from system cards, each per dollar of build and run cost. The parameters of the decision model are elicited as probability distributions from an ensemble of \voiNMembers{} LLMs, and Monte Carlo draws carry their uncertainty into rankings and group comparisons. We report rankings and comparisons rather than absolute values, because elicited parameters are imperfect: a bias the elicitors share across all evaluations moves both groups alike.
